\documentclass[11pt]{article}
\usepackage{amsmath}
\title{A successful clock notification and a historical clock bound are different evidence}
\begin{document}
\maketitle

A LinuxPTP status poll can show a new sync ingress with an old
offset estimate. The selected source offers a stronger fact:
a time-sync notification reached from the successful-update path.
That is still not evidence that a past acquisition used the same
offset. The missing step can be made explicit and checked when
measurement, rate, and delivery bounds are available.

\paragraph{Native notification.}
In the pinned \texttt{clock\_synchronize}, a failed offset update
returns before \texttt{NOTIFY\_TIME\_SYNC}, in both adjustment
and free-running branches. A successful update can notify even
if the servo reports unlocked. A successful jump can reset the
ingress field to zero before notification, so the event cannot
always date the offset. The same notification method selects
the \texttt{TIME\_STATUS\_NP} payload; that payload alone does
not identify whether it was polled or notified.

The C replay compiles the pinned synchronization, notification,
timestamp update/reset, and management-copy logic. Clock adjustments,
servo outcomes, units and transport are inert fixture inputs.
There is no actual clock or subscription. The controls include
failure with advanced ingress but no notification, a nonzero
ingress in success cases, jump with zero ingress, and later UTC,
runtime, serialization and subscription conditions under which
a fresh calculation does not yield an observed event. The
runtime and delivery cases test the supplied failure modes, not
a measured loss model. Thus a captured, source-identified event
with a matched nonzero ingress can distinguish the first two
cases. Absent delivery cannot distinguish estimate failure from
lack of a usable notification. Servo state, a time-base change,
and accuracy still require their own evidence. Native source
provenance is a premise; receipt of the same structure alone
does not establish it.

\paragraph{Exact historical bound under a declared rate law.}
Let the true offset be constant on acquisition interval \([u,v]\).
Outside it the same clock, with no unmodeled steps, has drift rate
bounded by \(\rho\ge0\) in absolute value. This means the offset
is Lipschitz continuous with constant \(\rho\), with no drift over
\([u,v]\). Let \(b(t)\) be the compressed time coordinate: \(t\)
before \(u\), \(u\) during \([u,v]\), and \(t-(v-u)\) after \(v\).
A valid trace factors through \(b\) as a \(\rho\)-Lipschitz
function on compressed time. To see necessity, split a
cross-epoch time difference at \(u\) and \(v\); the offset
changes by zero in the omitted part. Sufficiency follows because
\(b\) does not enlarge time differences.

Keep only authenticated measurement bounds \([a_i,c_i]\)
whose reports arrive before the decision, with each value dated
at its actual estimation time \(t_i\). Write \(s_i=b(t_i)\)
and include any prior interval at \(s=u\). The bounds are
consistent iff \(a_i-c_j\le\rho|s_i-s_j|\) for every \(i,j\).
Necessity follows from the drift law. For sufficiency the
function \(L(s)=\max_i(a_i-\rho |s-s_i|)\) is Lipschitz and
lies within every measured interval at its knot. The upper
dual \(U(s)=\min_i(c_i+\rho|s-s_i|)\) is also a feasible trace.
It follows that the exact permissible epoch offset interval
is \([L(u),U(u)]\): the two extremes are attained, and convex
interpolations realize all intermediate values.

For a query whether the old offset was at least \(3/4\),
answer yes iff \(L(u)\ge3/4\), no iff \(U(u)<3/4\), and otherwise
leave it open. Inconsistent evidence is a separate outcome.
This rule cannot use an accurate measurement delivered only
after the admission cut. Clock identity and no unmodeled steps
are necessary to apply the same \(\rho\), and path-delay bias
must be included in the input error intervals. They are not
deduced from the LinuxPTP notification.

\paragraph{Matched numerical check.}
With epoch \([0,4]\), prior \([0,1]\), drift bound \(1/10\),
and a current exact reading \(1/2\) at time 8, the exact historical
interval is \([1/10,9/10]\). It straddles the threshold even
though the current measurement is fresh. A previous estimate at
time \(-1\) in \([87/100,91/100]\) narrows it to
\([77/100,9/10]\) if delivered by time 8. Late delivery leaves
the original interval. The replay compares 320 exact interval
constructions against a separate finite-knot linear program.
Model times and errors are authored, not LinuxPTP calibration.
The result advances the interface check but does not establish a
historical query in a deployed system or full R4 acceptance.
\end{document}
